\documentclass[10pt]{article}

\usepackage[margin=0.5in]{geometry}
\usepackage{titlesec}
\titlespacing*{\section}{0pt}{0.5\baselineskip}{0.2\baselineskip}
\linespread{0.96}
\usepackage{amsmath, amssymb}
\usepackage{booktabs}
\usepackage{hyperref}
\usepackage{enumitem}
\setlist{nosep,leftmargin=*}
\setlength{\parskip}{0pt}
\setlength{\parindent}{0pt}

\begin{document}
\vspace*{-4em}
\begin{center}
{\Large \textbf{Project Proposal: Submodular Optimization and Greedy Approximation for Distributed Influence Maximization}}
\end{center}
\vspace*{-1em}

\section{Topic Selection}

We choose \textbf{Topic~3: Influence Maximization} from the course reference
list. Influence maximization asks how to select a limited number of seed nodes
in a social network so that information, behavior, products, or interventions
spread to as many nodes as possible. It is a fundamental model for viral
marketing, public-health interventions, and misinformation control, and it
provides a clean algorithmic connection between modern networked applications
and classical combinatorial optimization.

This project focuses on \textbf{submodular optimization and greedy
approximation for distributed influence maximization}. We study the problem
from two complementary perspectives:
\begin{itemize}
    \item \textbf{Algorithmic perspective:} How does the classical greedy
    algorithm exploit the submodularity of the influence function, and how can
    acceleration techniques such as lazy evaluation (CELF/CELF++) and
    reverse-influence sampling (TIM/IMM) reduce computation while preserving
    approximation guarantees?
    \item \textbf{Distributed perspective:} In large-scale networks, a single
    machine or agent may not have access to the entire graph. Each local agent
    observes a subgraph or estimates local diffusion effects, then communicates
    with other agents to approximate a globally effective seed set under a fixed
    budget. This setting naturally connects to distributed submodular
    maximization and multi-agent coordination problems.
\end{itemize}

\section{Problem Formulation}

Let $G=(V,E)$ be a directed or undirected graph with node set~$V$ and edge
set~$E$. Each edge $(u,v)$ carries a propagation probability
$p_{uv}\in[0,1]$. Given a budget~$k$, the goal is to choose a seed set
$S\subseteq V$ with $|S|\leq k$ that maximizes the expected number of
activated nodes:
\[
    \max_{S\subseteq V,\ |S|\leq k} \sigma(S),
\]
where $\sigma(S)$ is the expected influence spread from~$S$ under a stochastic
diffusion model.

\textbf{Diffusion models.}
We consider two classical models:
\begin{itemize}
    \item \textbf{Independent Cascade (IC):} When node~$u$ becomes active, it
    has a single chance to activate each inactive neighbor~$v$ with probability
    $p_{uv}$, independently.
    \item \textbf{Linear Threshold (LT):} Each node~$v$ draws a random
    threshold $\theta_v\in[0,1]$; node~$v$ becomes active when the sum of
    incoming edge weights from its active neighbors exceeds~$\theta_v$.
\end{itemize}

Under both models, $\sigma(S)$ is monotone and submodular~\cite{kempe2003},
meaning the marginal gain of adding a new node decreases as the seed set
grows:
\[
    \sigma(A\cup\{v\})-\sigma(A)
    \geq
    \sigma(B\cup\{v\})-\sigma(B),
    \quad A\subseteq B,\ v\notin B.
\]
Exact influence maximization is NP-hard. However, the classical greedy
algorithm of Nemhauser et al.~\cite{nemhauser1978} achieves a $(1-1/e)$
approximation guarantee for maximizing any monotone submodular function under
a cardinality constraint.

\section{Related Work}

\textbf{Kempe, Kleinberg, and Tardos}~\cite{kempe2003} introduced the
influence maximization problem for social networks, proved NP-hardness under
the IC and LT models, and showed that the greedy algorithm yields a
$(1-1/e)$-approximate solution. This is our \textbf{core foundation} for
connecting influence maximization with submodular greedy approximation.

\textbf{Nemhauser, Wolsey, and Fisher}~\cite{nemhauser1978} proved the
classical approximation guarantee for maximizing a nonnegative monotone
submodular function subject to a cardinality constraint. Their result provides
the theoretical foundation for all greedy-based influence maximization
algorithms.

\textbf{Leskovec et al.}~\cite{leskovec2007} proposed the Cost-Effective Lazy
Forward (CELF) optimization, which exploits the diminishing marginal gains of
submodular functions to skip unnecessary evaluations. CELF achieves up to
700$\times$ speedup over na\"ive greedy while preserving the same
approximation ratio. \textbf{Goyal, Lu, and
Lakshmanan}~\cite{goyal2011} further improved this with CELF++, achieving an
additional 35--55\% speedup.

\textbf{Tang, Xiao, and Shi}~\cite{tang2014} proposed TIM (Two-phase
Influence Maximization), which uses Reverse Influence Sampling (RIS) to
achieve near-optimal time complexity with a $(1-1/e-\varepsilon)$
approximation guarantee. \textbf{Tang, Shi, and Xiao}~\cite{tang2015}
subsequently proposed IMM (Influence Maximization via Martingales), which uses
martingale concentration inequalities to further tighten the sample complexity
and scale to networks with billions of edges.

\textbf{Mirzasoleiman et al.}~\cite{mirzasoleiman2013} proposed GreeDI, a
distributed protocol for submodular maximization under cardinality constraints.
GreeDI partitions the ground set across multiple machines, runs local greedy
on each partition, and merges the results. It achieves a constant-factor
approximation and provides the main basis for our distributed optimization
component.

\section{Algorithm and Technical Plan}

We plan to implement and systematically compare the following methods:

\textbf{Baseline methods:}
\begin{itemize}
    \item \textbf{Random:} select $k$ seed nodes uniformly at random.
    \item \textbf{Degree centrality:} select the $k$ nodes with highest
    (out-)degree as a simple graph-structural heuristic.
    \item \textbf{PageRank-based:} select the $k$ nodes with highest PageRank
    score.
\end{itemize}

\textbf{Greedy methods:}
\begin{itemize}
    \item \textbf{Standard greedy:} at each step, select the node with the
    largest estimated marginal influence gain via Monte Carlo simulation.
    \item \textbf{CELF / CELF++ lazy greedy:} maintain upper bounds on marginal
    gains and skip re-evaluation of candidates whose upper bound is below the
    current best, reducing the number of costly influence estimations.
\end{itemize}

\textbf{Scalable methods:}
\begin{itemize}
    \item \textbf{TIM / IMM:} generate random Reverse Reachable (RR) sets and
    select seeds by a weighted maximum coverage procedure. These methods
    achieve $(1-1/e-\varepsilon)$ approximation in near-linear time.
\end{itemize}

\textbf{Distributed optimization component:}
\begin{itemize}
    \item \textbf{GreeDI-style distributed greedy:} partition the graph across
    $m$ agents; each agent runs local greedy on its subgraph and returns its
    top-$k$ candidates; a central coordinator merges the $mk$ candidates and
    runs a final round of greedy selection. We will study how the number of
    partitions~$m$ and communication budget affect solution quality.
\end{itemize}

\textbf{Datasets.}
We will evaluate on both synthetic and real-world networks:
\begin{itemize}
    \item \textbf{Synthetic:} Erd\H{o}s--R\'enyi random graphs,
    Barab\'asi--Albert scale-free graphs, and Watts--Strogatz small-world
    graphs with varying sizes ($n \in \{500, 1000, 5000, 10000\}$).
    \item \textbf{Real-world:} NetHEPT (High-Energy Physics collaborations,
    ${\sim}15$K nodes, ${\sim}59$K edges) and Wiki-Vote (Wikipedia
    adminship votes, ${\sim}7$K nodes, ${\sim}104$K edges).
\end{itemize}

\textbf{Evaluation metrics:}
expected influence spread (estimated via Monte Carlo), wall-clock runtime,
number of influence evaluations, and scalability with respect to graph
size~$n$, seed budget~$k$, and number of Monte Carlo samples~$R$.

\section{Project Scope and Expected Goals}

\textbf{Core reproduction goals:}
\begin{itemize}
    \item Greedy selection achieves significantly larger influence spread than
    random, degree, and PageRank baselines, validating the effectiveness of
    submodular optimization.
    \item CELF/CELF++ preserves nearly the same solution quality as standard
    greedy while reducing runtime by one to two orders of magnitude.
    \item TIM/IMM achieves comparable influence spread with further runtime
    improvements, demonstrating the advantage of reverse-influence sampling.
    \item Influence spread increases with seed budget~$k$, while marginal
    gains decrease monotonically, confirming the submodularity of~$\sigma$.
\end{itemize}

\textbf{Distributed optimization goals:}
\begin{itemize}
    \item The GreeDI-style distributed scheme achieves a reasonable
    approximation ratio compared to centralized greedy.
    \item Solution quality degrades gracefully as the number of partitions~$m$
    increases, illustrating the tradeoff between communication cost and
    approximation quality in distributed submodular optimization.
\end{itemize}

\textbf{Optional further extensions:}
\begin{itemize}
    \item Compare the IC and LT diffusion models on the same graphs.
    \item Analyze the effect of Monte Carlo sample size on estimation variance.
    \item Implement a communication-limited consensus protocol where agents
    exchange partial information over multiple rounds.
    \item Discuss connections to distributed optimization with local
    information and global objectives.
\end{itemize}

\section{Team Information}

\begin{tabular}{@{}lll@{}}
\toprule
Member & Responsibility & Expected Contribution \\
\midrule
Yixin Chen & Modeling and theory & Problem formulation, related work, and complexity analysis \\
Bichi Zhang & Algorithms & Greedy, CELF/CELF++, TIM/IMM, GreeDI implementation and analysis \\
Yaoqin Ye & Experiments & Datasets, evaluation, plots, and distributed optimization analysis \\
\bottomrule
\end{tabular}

\begin{footnotesize}
\bibliographystyle{unsrt}
\bibliography{references}
\end{footnotesize}

\end{document}
